% Lösungen Einheit 1 von 4 – Terme aufstellen und berechnen (zusammenbau v0.1)
\einheitenkopf{Einheit 1 von 4 · Terme aufstellen und berechnen}

\erg{7}{a) $4x - 1$ \quad b) $3x + 8$ \quad c) $5x$ (auch $5 \cdot x$) \quad d) $4a + 3b$ \quad e) $x + 9 + x + 9 = 2x + 18$ \quad f) $2 \cdot (3x + 2)$ \quad g) Dreifaches $3x$, Differenz $3x - 7$, verdoppelt mit Klammer: $2 \cdot (3x - 7)$}
\erg{8}{a) $4 \cdot 5 - 3 = 17$}
\erg{9}{a) Zum Beispiel: Ein Kinobesuch kostet 12 € Eintritt, dazu kommen $x$ Getränke zu je 3 €; $x$ ist die Zahl der Getränke. Für $x = 2$: $12 + 3 \cdot 2 = 18$, der Besuch kostet 18 €.}
\erg{10}{a) Ben rechnet $4 + 2$ zuerst; Punkt vor Strich: $4 + 2 \cdot 3 = 4 + 6 = 10$. \quad b) $4 \cdot (x + 6)$ \quad c) $12 + 0{,}30n$; für 150 kWh: $12 + 0{,}30 \cdot 150 = 57$ €}
